\documentclass[aps,prl,twocolumn,showpacs,preprintnumbers]{revtex4-2}
\usepackage{amsmath,amssymb}
\usepackage{bm}

\newcommand{\SU}{\mathrm{SU}}
\newcommand{\Nc}{N_c}
\newcommand{\Nf}{N_f}
\newcommand{\MeV}{\,\mathrm{MeV}}
\newcommand{\GeV}{\,\mathrm{GeV}}
\newcommand{\TeV}{\,\mathrm{TeV}}

\begin{document}

\title{CKM and PMNS mixing from the Koide angle}

\author{A.\ Rivero}
\affiliation{Independent researcher}

\date{\today}

\begin{abstract}
We show that the Cartan angle $\delta = 2/9$ of the Koide
mass formula determines both the CKM and PMNS mixing matrices
with zero free parameters.
The Wolfenstein hierarchy $|V_{us}| = 2/9$,
$|V_{cb}| = 4/99$, $A = 9/11$ matches PDG data to
$0.9\%$, $4.3\%$, and~$1.0\%$.
The PMNS deviations from tribimaximal mixing are
$\theta_{12} = \pi/4 - \delta + \delta^2/N$,
$\theta_{23} = \pi/4 + \delta/N$,
$\theta_{13} = 2\delta/N = 4/27$~rad,
matching data to $0.7\%$, $0.08\%$, and~$0.9\%$.
All six mixing parameters are rational functions of $N = 3$
(the number of families), constituting a $1/N$ expansion
of $\SU(3)$ family symmetry.
Combining $\theta_C = \delta$ with the Gatto--Sartori--Tonin relation
gives the cross-sector prediction $m_d/m_s = \delta^2 = 4/81$ ($1.2\%$
from PDG).
The Koide formula $K = 2/3$ is an identity of
this expansion: $K = j(j{+}1)/N$ with $j = (N{-}1)/2 = 1$.
\end{abstract}

\maketitle

\section{Introduction}

The Koide mass formula~\cite{Koide:1983}
\begin{equation}
  K \equiv \frac{m_e + m_\mu + m_\tau}{(\sqrt{m_e}+\sqrt{m_\mu}+\sqrt{m_\tau})^2}
  = \frac{2}{3}
  \label{eq:koide}
\end{equation}
has held to one part in $10^5$ since 1981.
In the parametrisation $\sqrt{m_n} = \mu(1+\sqrt{2}\cos(\delta+2\pi n/3))$
of Brannen~\cite{Brannen:2006}, Foot's geometric
condition~$\theta = \pi/4$~\cite{Foot:1994} is equivalent to
Eq.~(\ref{eq:koide}), and the mass hierarchy is encoded in a
single Cartan angle~$\delta$.
Numerically $\delta = 0.22222 \pm 0.00001$~rad,
consistent with the exact value $\delta = 2/9$.

We show that $\delta$ determines the entire CKM and PMNS mixing
structure, extending the QLC of Raidal~\cite{Raidal:2004}
and the quark Koide phases of
Zenczykowski~\cite{Zenczykowski:2012}.


\section{Cabibbo angle = Koide angle}

Koide's original composite model~\cite{Koide:1983} predicted
\begin{equation}
  \tan\theta_C = \frac{\sqrt{3}(\sqrt{m_\mu}-\sqrt{m_e})}{2\sqrt{m_\tau}-\sqrt{m_\mu}-\sqrt{m_e}}\,.
  \label{eq:cabibbo}
\end{equation}
This expression is \emph{algebraically identical} to $\tan\delta$:
in the Koide parametrisation, the right-hand side of
Eq.~(\ref{eq:cabibbo}) equals $\tan\delta$ identically, for any
masses satisfying $K = 2/3$.
Thus $\theta_C = \delta$.
Numerically $\theta_C = 0.22223$~rad, matching the PDG value
$|V_{us}| = 0.2243$ to~$0.9\%$.


\section{CKM from the democratic resolvent}

The Wolfenstein parametrisation expresses the CKM matrix as
a power series in $\lambda = |V_{us}|$.
Setting $\lambda = \delta = 2/9$, the Wolfenstein~$A$ parameter
is determined by the \emph{resolvent} of the democratic matrix:

The democratic matrix $J = \mathbf{1}\mathbf{1}^T$ (all entries~1)
has eigenvalues $3$ (democratic direction) and $0$ (double,
breaking plane).
The resolvent of the zero-eigenvalue subspace at
$z = 1 + \delta$ gives
\begin{equation}
  A = \frac{1}{1+\delta} = \frac{N^2}{N^2+2} = \frac{9}{11}\,,
  \label{eq:A}
\end{equation}
where we used $\delta = 2/N^2$ with $N = 3$.

The resulting zero-parameter CKM magnitudes are:
\begin{equation}
  |V_{us}| = \frac{2}{9}\,,\qquad
  |V_{cb}| = \frac{4}{99}\,,\qquad
  A = \frac{9}{11}\,.
  \label{eq:ckm}
\end{equation}
Comparison with PDG~\cite{PDG:2024}:

\begin{center}
\begin{tabular}{lccc}
\hline
& Prediction & PDG & Off \\
\hline
$|V_{us}|$ & $2/9 = 0.2222$ & $0.2250 \pm 0.0007$ & $1.2\%$ \\
$|V_{cb}|$ & $4/99 = 0.0404$ & $0.0422 \pm 0.0008$ & $4.3\%$\footnotemark \\
$A$ & $9/11 = 0.818$ & $0.826 \pm 0.013$ & $1.0\%$ \\
\hline
\end{tabular}
\end{center}


\footnotetext{The PDG average uses the inclusive value.
The exclusive determination $|V_{cb}|_{\mathrm{excl}} = 0.0395 \pm 0.0008$
differs from the inclusive by $\sim 3\sigma$.
Our prediction $4/99 = 0.0404$ is only $1.1\sigma$ from the exclusive value.}

\section{Down-quark mass ratio}

Combining $\theta_C = \delta$ with the Gatto--Sartori--Tonin
relation~\cite{GST:1968} $|V_{us}| \approx \sqrt{m_d/m_s}$ gives
a zero-parameter cross-sector prediction:
\begin{equation}
  \frac{m_d}{m_s} = \delta^2 = \frac{4}{81} = 0.0494\,.
  \label{eq:md_ms}
\end{equation}
The PDG value $m_d/m_s = 0.0500 \pm 0.0044$ agrees to~$1.2\%$.
Together with the inverse Koide condition
$K^{-1}(d,s,b) = 2/3$ and the bion relation
$\sqrt{m_b} = 3\sqrt{m_s} + \sqrt{m_c}$,
the three down-type quark masses are determined from $m_c$ alone:
$m_s = 91.6$~MeV ($1.9\%$),
$m_b = 4141$~MeV ($1.0\%$),
$m_d = 4.52$~MeV ($3.2\%$).

\section{PMNS from the $1/N$ expansion}

The PMNS mixing angles are deviations from maximal mixing
($\pi/4$), controlled by~$\delta$:
\begin{align}
  \theta_{12} &= \frac{\pi}{4} - \delta + \frac{\delta^2}{N}
  = \frac{\pi}{4} - \frac{2}{9} + \frac{4}{243}\,,
  \label{eq:t12} \\
  \theta_{23} &= \frac{\pi}{4} + \frac{\delta}{N}
  = \frac{\pi}{4} + \frac{2}{27}\,,
  \label{eq:t23} \\
  \theta_{13} &= \frac{2j(j{+}1)}{N^3} = \frac{4}{27}\;\mathrm{rad}\,,
  \label{eq:t13}
\end{align}
where $j = (N{-}1)/2 = 1$ and $N = 3$.

\begin{center}
\begin{tabular}{lccc}
\hline
& Prediction & NuFIT~5.2 & Off \\
\hline
$\sin^2\theta_{12}$ & $0.300$ & $0.304 \pm 0.012$ & $1.3\%$ \\
$\sin^2\theta_{23}$ & $0.574$ & $0.573 \pm 0.018$ & $0.1\%$ \\
$\sin^2\theta_{13}$ & $0.0218$ & $0.0222 \pm 0.0007$ & $1.9\%$ \\
\hline
\end{tabular}
\end{center}

The PMNS correction coefficients $(-1, +1/3, +2/3)$ are
\emph{traceless} (they sum to zero) and \emph{unique}:
among all traceless triples with denominators $\leq 3$,
this is the only one giving all three angles within $5\%$
of data. Moreover, scanning~$N = 2,\ldots,10$ in
Eqs.~(\ref{eq:t12})--(\ref{eq:t13}), only $N = 3$ gives
$\chi^2 < 1$ ($\chi^2 = 0.5$); $N = 2$ gives $604$, $N = 4$
gives~$218$.


\section{The $1/N$ expansion}

All parameters are rational functions of $N$:
\begin{center}
\begin{tabular}{llll}
\hline
Parameter & Formula & Value & Order \\
\hline
$\delta$ & $2/N^2$ & $2/9$ & $O(1/N^2)$ \\
$K$ & $j(j{+}1)/N$ & $2/3$ & $O(1/N)$ \\
$|V_{us}|$ & $2/N^2$ & $2/9$ & $O(1/N^2)$ \\
$A$ & $N^2/(N^2{+}2)$ & $9/11$ & $1{-}O(1/N^2)$ \\
$\theta_{12}{-}\pi/4$ & $-2/N^2{+}4/N^4$ & & $O(1/N^2)$ \\
$\theta_{23}{-}\pi/4$ & $2/N^3$ & $2/27$ & $O(1/N^3)$ \\
$\theta_{13}$ & $4/N^3$ & $4/27$ & $O(1/N^3)$ \\
\hline
\end{tabular}
\end{center}

The lepton mass hierarchy, the CKM matrix, and the PMNS matrix
are all the $1/N$ expansion of $\SU(N)$ family symmetry at~$N = 3$.


\section{Statistical assessment}

With $|V_{us}|$ treated as input (determining~$\delta$),
the remaining five predictions have
$\chi^2/\mathrm{dof} = 6.2/5 = 1.25$ ($p = 0.28$).
The pulls are: $|V_{cb}|$ $(-2.3\sigma)$,
$A$ $(-0.6\sigma)$,
$\sin^2\theta_{12}$ $(-0.3\sigma)$,
$\sin^2\theta_{23}$ $(+0.04\sigma)$,
$\sin^2\theta_{13}$ $(-0.6\sigma)$.
The fit is acceptable at the 28\% level.

If instead $\delta = 2/9$ is used without adjustment,
$|V_{us}|$ contributes $-4.2\sigma$, reflecting the $1.2\%$
discrepancy between $2/9 = 0.2222$ and the PDG value $0.2250$.
This tension may be resolved by the instanton-generated
potential $V(\delta) = -\cos(3\delta) + (1/\pi)\cos(6\delta)$,
which shifts the minimum to $\delta = 0.2225$.

\section{Discussion}

Six mixing parameters (three CKM magnitudes, three PMNS angles)
follow from a single input: $N = 3$.
The physical origin is a confining $\SU(3)_F$ family gauge
theory~\cite{Rivero:2024} whose Koide formula is a Casimir
identity~\cite{CasimirKoide}: $K = \langle J_3^2\rangle$
iff $N(N^2{-}1) = 24$ iff $N = 3$.

The critical experimental test is the solar angle.
Our prediction $\sin^2\theta_{12} = 0.300$ is $1.3\%$
below the NuFIT central value $0.304$.
JUNO~\cite{JUNO} will measure $\sin^2\theta_{12}$ to~$\pm 0.003$,
testing this prediction at~$\sim 1.3\sigma$.
A measurement confirming $\sin^2\theta_{12} < 0.295$ would
strongly support the $1/N$ expansion.

The CKM--PMNS hierarchy is structural:
CKM elements scale as $\delta^{|i-j|}$ (a power of $1/N^2$),
while PMNS corrections scale as $1/N^{|i-j|-1}$.
PMNS is near-maximal because its corrections are $O(1/N)$;
CKM is hierarchical because its elements are $O(1/N^2)$.
Both originate from the same $\SU(3)_F$, but at different orders
in the $1/N$ expansion.

Two quantities remain unpredicted: the CKM and PMNS CP-violating
phases.
A candidate mechanism for the Cartan angle~$\delta$ is the
minimum of the angular potential
$V(\delta) = -\cos(3\delta) + (1/\pi)\cos(6\delta)$,
which gives $\delta = 2/9$ to~$0.12\%$.

Assuming a neutrino Koide angle~$\delta_\nu \approx 0.47$~rad
(fitted to mass-squared differences), the framework predicts
normal ordering with $m_1 \approx 0.35$~meV, $m_3 \approx 51$~meV,
$\sum m_\nu \approx 60$~meV (below the Planck bound), and
effective Majorana mass $m_{ee} \approx 4$~meV (within
reach of tonne-scale $0\nu\beta\beta$ experiments).

\begin{acknowledgments}
The author thanks M.~Porter for identifying Seiberg duality
as the mechanism for the sBootstrap during extended discussions
(2011--2024), C.~Brannen for the circulant matrix parametrisation
of the Koide formula, and the late M.~Sheppeard (Kea) for
early contributions connecting Koide to representation theory.
During the preparation of this work the author used Claude
(Anthropic) for algebraic verification, numerical checks,
and drafting assistance.
The author reviewed and edited all content and takes full
responsibility for the publication.
\end{acknowledgments}

\begin{thebibliography}{99}

\bibitem{Koide:1983}
  Y.~Koide,
  Phys.\ Lett.\ B \textbf{120}, 161 (1983).

\bibitem{Brannen:2006}
  C.A.~Brannen,
  \texttt{hep-ph/0606057} (2006).

\bibitem{Foot:1994}
  R.~Foot,
  arXiv:hep-ph/9402242 (1994).

\bibitem{PDG:2024}
  R.L.~Workman \textit{et al.} (Particle Data Group),
  Prog.\ Theor.\ Exp.\ Phys.\ \textbf{2022}, 083C01 (2022)
  and 2024 update.

\bibitem{Rivero:2024}
  A.~Rivero,
  Eur.\ Phys.\ J.\ C \textbf{84}, 1058 (2024).

\bibitem{CasimirKoide}
  A.~Rivero, in preparation.

\bibitem{Raidal:2004}
  M.~Raidal,
  Phys.\ Rev.\ Lett.\ \textbf{93}, 161801 (2004),
  arXiv:hep-ph/0404046.

\bibitem{Zenczykowski:2012}
  P.~Zenczykowski,
  Phys.\ Rev.\ D \textbf{86}, 117303 (2012),
  arXiv:1210.4125.

\bibitem{GST:1968}
  R.~Gatto, G.~Sartori, and M.~Tonin,
  Phys.\ Lett.\ B \textbf{28}, 128 (1968).

\bibitem{JUNO}
  F.~An \textit{et al.} (JUNO Collaboration),
  J.\ Phys.\ G \textbf{43}, 030401 (2016).

\end{thebibliography}

\end{document}
